\documentclass[11pt]{article}
\usepackage{mathunal-base}
\mathunalfuente{serif}
\usepackage{amssymb}
\usepackage{graphicx}
\graphicspath{{fig/}}

\providecommand{\nota}[1]{\par\smallskip\noindent{\small\itshape\color{unalblue}#1}\par\smallskip}
\newcommand{\resp}{\underline{\hspace{3cm}}}
\newcommand{\vf}{\underline{\hspace{1.2cm}}}

\begin{document}

{\large\bfseries Segundo Parcial (Supletorio) de Matemáticas Discretas --- 6 de junio de 2022}

\medskip
\noindent Escuela de Matemáticas. Universidad Nacional de Colombia, Sede Medellín.

\medskip
\noindent\textit{Duración: 1 hora y 50 minutos. Puntaje (sólo para uso oficial): Ej.\ 1--2--3 (40), Ej.\ 4 (20), Ej.\ 5 (20), Ej.\ 6 (20). Solucione las preguntas en los espacios en blanco asignados a cada una. No está permitido sacar hojas en blanco ni ningún tipo de apuntes durante el examen. En los ejercicios 4 a 6 se tiene en cuenta el procedimiento.}

\nota{Transcripción: los diagramas del Problema 1 se reproducen a partir de las figuras digitales del examen original.}

\begin{enumerate}
\item[]\textbf{Problema 1.} (16 puntos) Para cada una de las siguientes figuras, escriba V si ésta se puede dibujar sin levantar la mano del papel ni repetir trazos, o F en caso contrario.

\begin{center}
\begin{tabular}{cccc}
\includegraphics[height=2.6cm]{supl_f1abc_0} &
\includegraphics[height=2.6cm]{supl_f1abc_1} &
\includegraphics[height=2.6cm]{supl_f1abc_2} &
\includegraphics[height=2.6cm]{supl_f1d_0} \\[2pt]
(a) (4 puntos) & (b) (4 puntos) & (c) (4 puntos) & (d) (4 puntos) \\
Respuesta: \vf & Respuesta: \vf & Respuesta: \vf & Respuesta: \vf
\end{tabular}
\end{center}

\item[]\textbf{Problema 2.} (12 puntos) Considere el grafo $G$ cuyos vértices son los números $\{2, 3, 4, 5, 6, 7, 8, 9, 10\}$. Dos vértices del grafo $G$ están unidos por un lado si y sólo si los correspondientes números son relativamente primos. Calcule el número cromático de $G$.

\vspace{2pt}
Respuesta: \resp

\item[]\textbf{Problema 3.} (12 puntos) Para cada literal, escriba V si es posible construir un grafo con las propiedades mencionadas y F en caso contrario.
\begin{enumerate}[label=(\alph*)]
\item (3 puntos) Un grafo con $7$ vértices en el que todos los vértices tienen grado $3$. \hfill Respuesta: \vf
\item (3 puntos) Un grafo con $6$ vértices, número cromático $2$ y $10$ aristas. \hfill Respuesta: \vf
\item (3 puntos) Un grafo planar con $4$ vértices, $6$ aristas y $4$ caras. \hfill Respuesta: \vf
\item (3 puntos) Un grafo planar con $8$ vértices en que todos los vértices tienen grado $6$. \hfill Respuesta: \vf
\end{enumerate}

\item[]\textbf{Problema 4.} (15 puntos) Encuentre números enteros $s$ y $t$ tales que
\[
1 = 1247\,s + 245\,t.
\]
Respuesta: $s = \underline{\hspace{2cm}}$, \quad $t = \underline{\hspace{2cm}}$.

\item[]\textbf{Problema 5.} (15 puntos) Encuentre un número natural $u$ tal que
\[
257\,u \equiv 7 \pmod{333}.
\]
Respuesta: \resp

\item[]\textbf{Problema 6.} (20 puntos) Alicia quiere recibir un mensaje seguro usando el protocolo RSA. Ella publica su llave pública $(N = 323,\ e = 5)$. Berta quiere mandarle el mensaje $M$ a Alicia. Para eso lo encripta usando la llave pública, y publica el mensaje encriptado:
\[
\text{Encriptado} = 122.
\]
El mensaje original $M$ es: \resp

\end{enumerate}

\vfill
\noindent{\small\itshape Transcripción digital --- MathUNAL}

\end{document}
